
\textbf{Semantic Representations and Emergent Invariance in Language Models.}
Prior work shows that transformer-based language models develop structured representations across depth, with lower layers encoding lexical and syntactic information and higher layers capturing contextual and semantic properties~\cite{ethayarajh-2019-contextual,tenney-etal-2019-bert}. Probing studies further demonstrate that linguistic structures such as syntax and dependency relations can be recovered from these representations~\cite{hewitt-manning-2019,rogers-etal-2020}. Geometric analyses of BERT additionally suggest that semantic and syntactic information occupy structured subspaces in latent space, with word senses exhibiting fine-grained organization~\cite{reif2019visualizing}. Together, these findings indicate that language models encode semantic structure beyond surface form. 
This line of work also implies an empirical form of semantic invariance: inputs that differ in wording but preserve meaning often yield similar representations. Analyses of contextualized embeddings and their reductions to static representations further show that semantic information remains stable across contexts~\cite{ethayarajh-2019-contextual,bommasani-etal-2020-interpreting}.

Existing studies primarily establish that semantic structure is present in language-model representations. However, a local geometric characterization of how semantic invariance emerges across layers, how paraphrastic variation separates from semantic variation, and how such invariant structure can be operationally identified, remains less explored.

\textbf{Representation Geometry and Subspace Analysis in Neural Networks.}
A central line of work studies neural representations through geometric and statistical comparisons of activation spaces. Representational similarity analysis (RSA) and related methods quantify relationships across inputs, layers, and models. Widely used techniques include canonical correlation analysis (CCA) and its variants such as singular vector CCA (SVCCA)~\cite{raghu2017svcca} and projection-weighted CCA (PWCCA)~\cite{morcos2018pwcca}, as well as centered kernel alignment (CKA), which is invariant to orthogonal transformations and isotropic scaling~\cite{kornblith2019similarity}. These methods enable systematic comparison of representations and are widely used to analyze training dynamics, layer correspondence, and cross-model similarity. 
Complementary approaches examine the structure of representation spaces directly. Linear probes show that linguistic and semantic features can be recovered from learned representations, while geometric analyses indicate that such information resides in structured subspaces~\cite{hewitt-manning-2019,reif2019visualizing}. Together, these results suggest that neural representations admit meaningful low-dimensional structure that can be characterized through linear projections and subspace decompositions.

Despite these advances, existing approaches mainly measure similarity or identify directions associated with specific properties. They do not explicitly provide mechanisms for separating semantic-preserving from semantic-changing variation. Consequently, while prior work reveals structure in representation space, a principled framework for isolating invariant subspaces that distinguish semantic content from nuisance variation remains limited. This motivates our formulation of invariant subspace discovery as a contrastive geometric decomposition, where invariant directions arise from a generalized eigenvalue operator that separates semantic and nuisance components.

\textbf{Model-Specific Representations and Attribution via Internal Activations.}
A growing body of work suggests that language models encode model-specific structure despite similar external behavior. Evidence arises from transfer studies, where soft prompts optimized for one model often fail to generalize to others, indicating interaction with model-specific latent pathways~\cite{lester-etal-2021-power}. Related observations come from adversarial trigger analysis, where input-agnostic triggers reveal global biases and model-dependent vulnerabilities~\cite{wallace-etal-2019-universal}. These results suggest that similar linguistic functions may be implemented through different internal geometries. 
Model attribution has also been studied via generated text, watermarking, and explicit identification signals. Prior work explores source-model identification from outputs~\cite{antoun-etal-2024-text}, watermark-based attribution~\cite{lu-etal-2024-source, dasgupta2025watermarking}, and instruction-based identification through lightweight training interventions~\cite{xu-etal-2024-instructional}.

While these approaches demonstrate that model identity can be detected, they typically rely on output-level artifacts, task-specific classifiers, watermarking mechanisms, or training-induced signatures. In contrast, our work uses invariant latent structure itself as a compact semantic signature. Attribution is therefore not treated as a standalone identification mechanism, but as a downstream validation that invariant subspaces preserve both semantic stability and model-specific representation geometry.













